\chapter{Introduction}

\section{What this is}

This is a particle physics engine written in CUDA. It solves two problems that
sit at opposite ends of the memory access spectrum, which is the reason both
are here.

The first is direct all pairs gravitational N-body dynamics. Every particle
attracts every other, so the force evaluation is $O(N^2)$ arithmetic over $O(N)$
data. Access is perfectly regular and entirely predictable.

The second is smoothed particle hydrodynamics, in which a particle interacts
only with the neighbours inside a finite smoothing radius. The work is nearly
linear in the particle count, but which particles are neighbours changes every
step and has to be discovered before any physics can be evaluated. Access is
irregular and data dependent.

A GPU rewards the first pattern and punishes the second, and the techniques
that matter are different in each case. Putting both in one engine, behind one
configuration format and one integrator, makes the contrast concrete.

\section{The two commitments}

\paragraph{No number without a run.} Every performance figure in this report
was produced by a program in this repository, recorded to a comma separated
file with the commit that produced it, and read back into this document
mechanically. The only typed values are \texttt{v1.0.0} values in correction records and
a few figures from one off runs that no file records, such as the double precision
truncation error in Section~\ref{sec:gates}. Where I had an expectation
before measuring, the expectation is stated and then compared with what
actually happened. Where a published claim turned out to be wrong,
Chapter~\ref{ch:corrections} says so.

\paragraph{Correctness is enforced, not observed.} A particle simulation that
is subtly wrong still looks like a particle simulation. Fluid still sloshes and
stars still orbit when the force law has a sign error in one term or the
neighbour search quietly counts some particles twice. The defence is a set of
conservation and distribution tests with numerical thresholds, run by the test suite on
a machine with a GPU, that fail loudly. There are \GateCount{} of them and
Chapter~\ref{ch:results} reports what each one measured.

These two commitments interact. The benchmark builds use \texttt{fast math},
which lowers the reciprocal square root to a hardware approximation. The physics gates
should not depend on a flag like that. So the build compiles the device library twice from
identical sources, once with fast math for the benchmarks and once without for
the gates. A gate that a compiler flag can relax is not a gate.

\section{The reference machine}

Everything here was measured on one machine: an \GpuName{} of compute
capability \GpuCompute{} with \GpuSms{} streaming multiprocessors on a
\GpuBusBits{} bit memory bus, driven from Ubuntu under the Windows Subsystem
for Linux with CUDA \CudaRuntime{}. Its theoretical peak memory bandwidth is
\GpuPeakBandwidth{} GB/s. A streaming copy kernel achieves
\GpuMeasuredBandwidth{} GB/s, or \GpuBandwidthEfficiency{} percent of that
peak, and it is the measured figure rather than the theoretical one that every
bandwidth statement in this report is compared against. Quoting utilisation
against a peak no kernel can reach flatters every kernel by the same factor
while saying nothing about any of them.

\section{How to read this report}

Chapter~\ref{ch:background} covers the GPU memory hierarchy and the physics of
both solvers, including why velocity Verlet rather than anything else.
Chapter~\ref{ch:methodology} describes the kernel ladder, the timing protocol
and the design of the validation suite. Chapter~\ref{ch:implementation}
describes what was built. Chapter~\ref{ch:results} is the measurements.
Chapter~\ref{ch:discussion} interprets them, including the results that
contradicted my expectations. Chapter~\ref{ch:corrections} lists what
\texttt{v1.0.0} got wrong. A companion document,
\texttt{debug\_report.pdf}, records the defects found along the way and how
each was diagnosed.
